\documentclass[11pt,a4paper]{article}
\usepackage[a4paper,margin=2.2cm]{geometry}
\usepackage[T1]{fontenc}
\usepackage{lmodern}
\usepackage{amsmath,amssymb}
\usepackage{array,booktabs}
\usepackage{enumitem}
\usepackage{xcolor}
\usepackage{fancyhdr}
\usepackage{lastpage}
\pagestyle{fancy}
\fancyhf{}
\cfoot{Page \thepage\ of \pageref{LastPage}}
\setlength{\parindent}{0pt}
\setlength{\parskip}{5pt}
\newcounter{qnum}
\newenvironment{examq}[2]{\stepcounter{qnum}\par\noindent\begin{minipage}{\linewidth}\setlength{\parskip}{5pt}%
{\large\bfseries Question \theqnum: #1}\hfill{\bfseries(#2 marks)}\par\smallskip}{\end{minipage}\par\bigskip}
\newcommand{\qpart}[2]{\par\textbf{(#1)} \textit{[#2]}\ }
\lhead{Decision Theory (DETH) --- Practice quiz}
\rhead{Name: \rule{5.5cm}{0.4pt}}
\begin{document}
\begin{center}
{\LARGE\bfseries Decision Theory (DETH)}\\[3pt]
{\Large Practice quiz --- MDPs, Bellman equations, policy and value iteration}\\[6pt]
Iman Mossavat $\cdot$ Fontys ICT $\cdot$ 26 November 2026\\
45 minutes $\cdot$ 5 questions $\cdot$ 27 marks
\end{center}
\textbf{Student name:} \rule{9cm}{0.4pt}

\fbox{\parbox{\dimexpr\linewidth-2\fboxsep-2\fboxrule}{\textbf{Marks are given for sound reasoning.} Show the formula, the numbers you put in, and the result.
A correct final number on its own earns nothing. Write your answers on the answer sheet.}}

Allowed: a simple calculator. No laptop, no phone, no gen-AI. Round to \textbf{two decimals}.

\subsection*{The problem (used in all questions): the coffee cart}
A coffee cart decides at the start of every hour what to do. Its stock is \textbf{Stocked} or \textbf{Low}.
With a stocked cart it sells at the \emph{normal price} or runs a \emph{happy hour}, which earns more now but empties the cart faster.
With a low stock it can \emph{keep selling} (a supplier sometimes drops by) or \emph{restock}, which costs money.
Decisions, including restocking, happen only at the start of an hour, never in the middle of one.
So a cart with low stock that keeps selling often sells out early in the hour and then has to wait: the price is the same, but it sells far fewer cups.
The reward is the profit of the hour. The cart runs with no fixed closing time.

\begin{center}\renewcommand{\arraystretch}{1.3}
\begin{tabular}{llccc}
\toprule
State $s$ & Action $a$ & Reward $r$ & $P(\text{next}=\text{Stocked})$ & $P(\text{next}=\text{Low})$ \\ \midrule
Stocked & Normal price & $10$ & $0.7$ & $0.3$ \\
Stocked & Happy hour & $14$ & $0.1$ & \textbf{?} \\
Low & Keep selling & $1$ & $0.2$ & $0.8$ \\
Low & Restock & $-3$ & $1$ & $0$ \\
\bottomrule
\end{tabular}
\end{center}
\textbf{The discount factor is $\boldsymbol{\gamma = 0.8}$} (used in all questions).

\begin{examq}{The model}{6}
\qpart{a}{1 mark} Find the missing probability. Which rule did you use?
\qpart{b}{3 marks} On sunny days the cart sells much faster, and sunny weather usually lasts for several hours.
The weather is not part of the state.
Is this model Markov? Explain. How would you change the state, and how many states do you then have?
\qpart{c}{2 marks} Is this task episodic or continuing? Explain. Now imagine that every single step, forever, gives the largest
reward in the table. What is the return $G_t$ then?
\end{examq}

\begin{examq}{Return and expectation}{4}
\qpart{a}{2 marks} In one run the first three rewards are $R_1=14$, $R_2=1$, $R_3=-3$,
and the return from time 3 is $G_3=20$.
Compute the return $G_0$ with $\gamma=0.8$.
\qpart{b}{2 marks} Start in Stocked and choose \emph{Happy hour}. In the next state choose \emph{Normal price} if it is Stocked and
\emph{Keep selling} if it is Low. Draw the tree and compute the expected value of $R_1+\gamma R_2$.
\end{examq}

\begin{examq}{Evaluating a policy}{7}
The policy $\pi_0$ chooses \emph{Normal price} in Stocked and \emph{Keep selling} in Low.
\qpart{a}{3 marks} Write the Bellman equation for $v_{\pi_0}(\text{Stocked})$ and for $v_{\pi_0}(\text{Low})$.
\qpart{b}{2 marks} Solve the two equations.
\qpart{c}{2 marks} Compute $q_{\pi_0}(\text{Stocked},\text{Happy hour})$ and $q_{\pi_0}(\text{Low},\text{Restock})$.
\end{examq}

\begin{examq}{Policy iteration}{4}
\qpart{a}{2 marks} Use your results of Question 3 to do one policy improvement step. Give the new policy $\pi_1$ and say why.
\qpart{b}{2 marks} Can you already conclude that $\pi_1$ is optimal? What does policy iteration do next, and when does it stop?
\end{examq}

\begin{examq}{Value iteration}{6}
Start with $V_0(\text{Stocked})=V_0(\text{Low})=0$.
\qpart{a}{2 marks} Compute $V_1$ from $V_0$ for both states. Which action gives the maximum in each state?
\qpart{b}{3 marks} Compute $V_2$ from $V_1$ for both states. Which action gives the maximum in each state?
\qpart{c}{1 mark} Explain in words why a best action can differ between the first and the second sweep.
\end{examq}
\end{document}
